\documentclass[a4paper,11pt]{scrartcl}

\usepackage[utf8]{inputenc}
\usepackage[ngerman]{babel}
\usepackage[T1]{fontenc}
\usepackage{amsmath}
\usepackage[left=2cm, right=2cm, top= 2cm, bottom = 2cm]{geometry}
\usepackage{setspace}
\usepackage{times}
\usepackage[svgnames]{xcolor}
\setcounter{tocdepth}{3}
\usepackage{graphicx}
\usepackage{tocstyle}
\usetocstyle{KOMAlike}
\usepackage{multicol}
\usepackage{url}
\usepackage{epstopdf}
\usepackage{verbatim}
\usepackage{listings} 


	
\usepackage{pdfpages}
\usepackage{capt-of}
\usepackage{comment}
%\pagestyle[headings] %normale vordefinierte Kopfzeile. Meist der Titel des auf dieser Seite gezeigten Inhaltsabschnitts
\usepackage{fancyhdr}
\bibliographystyle{unsrt} 

\pagestyle{fancy}
%Selbsteingestellte Kopf- und Fußzeilenbefehle
\lhead{\leftmark}
\chead{}
\rhead{Reverse Innovation} %kopfzeile, rechts-center-links

\rfoot{}
\cfoot{\thepage}
\lfoot{}

%Befehle um Kopfund Fuszeile Anzupassen:
%\thepage - Seitennummer
%\leftmark - Aktueller Kapitelname "Kapitel X. Name des Kapitels"
%\rightmark - Aktueller Untertitelname "X.X Name des Kapitels
%\chaptername - Das Wort "Kapitel"
%\thechapter - Aktuelle Kapitelnummerierung
%\thesection - Aktuelle Unterkapitelnummerierung
%\slshape - Kursiv und Sanseriv

\usepackage[left=2cm,right=2cm,top=2cm,bottom=2cm]{geometry}
\author{Simon  Stadlinger}
\begin{document}
%_________________Titelseite_______________________________________
	\begin{titlepage}
		\centering
		\includegraphics[scale=.6]{hpi_logo.jpg}\\\vspace{2cm}
		{\LARGE Modul "`Unternehmerisches Denken und Handeln"'}\\\vspace{.4cm}
		{\Huge \textbf{Reverse Innovation}}\\\vspace{2cm}
		
	\end{titlepage}
\newpage
\thispagestyle{empty}
\section*{}
\newpage
\pagestyle{fancy}
%________________Hauptteil_________________________________________
%================Inhaltsverzeichnis================================
\tableofcontents\newpage
%================Einleitung========================================
\section{Von der Globalisierung zu Reverse Innovation}
ist ne schwierige Story
\subsection{Motivation}
ist ne lange Geschichte
%________________Quellen___________________________________________

\end{document}